\section{Order optimization and collision-path phases}
\label{sec:v26-phases}
The risk profile and its logarithmic inverse admit optimization on the
same subset lattice. This gives a finite procedure for computing the
comparison profile, rather than the exact minimax distortion constant.

\begin{proposition}[Subset recursion]\label{prop:v26-dynamic}
For $S\subseteq V$ put $h(S)=\Theta_{\delta(S)}(M,a)$, with
$h(\varnothing)=h(V)=0$. Define $f(\varnothing)=0$ and, in increasing
cardinality,
\begin{equation}\label{eq:v26-dynamic}
 f(S)=\max\left\{h(S),\min_{x\in S}f(S\setminus\{x\})\right\}.
\end{equation}
Then $f(V)=Q_G(M,a)$, and minimizing predecessors give an optimal order.
Replacing $h(S)$ by $\sum_{e\in\delta(S)}b_e(\varepsilon,a)$ gives
$f(V)=W_G(\varepsilon,a)$. Once the cut costs are supplied, either
recursion uses $O(v2^v)$ comparisons and $O(2^v)$ stored values.
\end{proposition}
\begin{proof}
A path ending at $S$ has a last predecessor $S\setminus\{x\}$.
Its largest cost is the larger of $h(S)$ and the largest cost on its
predecessor path. Minimize first among paths with this predecessor,
and then among $x\in S$. The identity
$\min_x\max\{c,a_x\}=\max\{c,\min_xa_x\}$ proves
\eqref{eq:v26-dynamic} by induction on $|S|$. Backtracking realizes
each minimizing choice. There are $2^v$ subsets and
$\sum_S|S|=v2^{v-1}$ predecessor comparisons. The bit recursion is
the identical argument with its additive cut costs. These are
arithmetic/comparison counts for supplied real costs, not a claim
that arbitrary prior moments have a finite exact compiler or that
the dependence on $v$ is polynomial.
\end{proof}

\subsection{Analytic paths}
Let $a(\theta)$, $0\le\theta\le\theta_0$, be a real-analytic path
in the admitted compact calibration chamber, analytic in a
neighbourhood of zero. For an edge $e$, let $\tau_{e,i}(\theta)$
be its normalized positive formal future nodes. Retain their formal
labels even when two node functions coincide identically. For each
pair set
\[
 \nu_{e,ij}=\operatorname{ord}_{\theta=0}
                  (\tau_{e,i}(\theta)-\tau_{e,j}(\theta))
                  \in\N\cup\{0,\infty\}.
\]
Order $\infty$ means an identically zero difference near zero. Put
$\eta_{e,0}=0$, and, for $1\le\ell\le p_e$,
\begin{equation}\label{eq:v26-contact-volumes}
 \eta_{e,\ell}=\min_{|I|=\ell}
                    \sum_{\{i,j\}\subset I}\nu_{e,ij}.
\end{equation}
A singleton sum is zero. An infinite term is omitted in subsequent
maxima. For $s\ge0$ define
\begin{equation}\label{eq:v26-tropical-width}
 \begin{split}
 w_e(s)&=\max_{0\le\ell\le p_e}\{\ell s-\eta_{e,\ell}\},\\
 W_\pi(s)&=\max_{1\le j<v}\sum_{e\in\delta(S_j^\pi)}w_e(s),
 \qquad W(s)=\min_\pi W_\pi(s).
 \end{split}
\end{equation}
Here $w_e$ is a leading bit coefficient, not a finite-$\theta$ cost.

\begin{theorem}[Finite collision-path phase law]\label{thm:v26-phases}
On every compact interval $I\subset(0,\infty)$, uniformly in $s\in I$,
\begin{equation}\label{eq:v26-phase-law}
 B_G^{*,\mathrm{av}}(\theta^{2s},a(\theta))
 =W(s)\log_2(1/\theta)+O(1),\qquad
 B_G^{*,\mathrm{max}}(\theta^{2s},a(\theta))
 =W(s)\log_2(1/\theta)+O(1).
\end{equation}
The function $W$ is continuous, nondecreasing and piecewise affine,
with finitely many pieces, nonnegative integer slopes and rational
breakpoints. There is a finite subdivision on whose open intervals
one deterministic order realizes $W$ throughout. Convexity of $W$
is not required and need not hold.

For either risk criterion let $B_\pi^*$ denote the optimal bit
requirement when the order is fixed to $\pi$. Then
\begin{equation}\label{eq:v26-regret-law}
 \min_\pi\sup_{s\in I}
 \{B_\pi^*(\theta^{2s},a(\theta))-B_G^*(\theta^{2s},a(\theta))\}
 =\mathfrak g_I\log_2(1/\theta)+O(1),
\end{equation}
where
$\mathfrak g_I=\min_\pi\sup_{s\in I}\{W_\pi(s)-W(s)\}$.
For a closed interval $I$, its endpoints and the breakpoints of these
finitely many functions suffice to compute every supremum.
\end{theorem}
\begin{proof}
A nonzero analytic node difference has a finite integral order and a
nonzero leading coefficient. On a sufficiently short positive interval
its absolute value is bounded above and below by positive multiples
of $\theta^{\nu_{e,ij}}$. A determinant product indexed by $I$
therefore has order $\sum_{\{i,j\}\subset I}\nu_{e,ij}$, or is
identically zero. There are finitely many such products, and their
maximum has the smallest of their finite orders. Thus
\[
 V_{e,\ell}(a(\theta))\asymp\theta^{\eta_{e,\ell}}
\]
for finite $\eta_{e,\ell}$; for infinite order the volume is zero.
All comparisons are simultaneous over the finite edge and index sets.
Writing $U=\log_2(1/\theta)$ and using the zero index, the bit formula
\eqref{eq:v25-edge-bits} gives
$b_e(\theta^{2s},a(\theta))=Uw_e(s)+O(1)$, with an error independent
of $s$. Finite sums, maxima and minima change by at most the largest
accumulated error. Corollary~\ref{cor:v25-graph-bits} proves
\eqref{eq:v26-phase-law}; its fixed-order version from
Proposition~\ref{prop:v25-fixed-order} proves
$B_\pi^*=UW_\pi+O(1)$ simultaneously for every $\pi$ and $s$.

Each $w_e$ is a maximum of finitely many affine functions with integral
slope and intercept. On each common interval of affinity, every cut
sum is affine with integral coefficients. Subdivide further at every
intersection of these finitely many affine functions. Their relative
orders are fixed on an open interval of the resulting subdivision,
so all cut maxima and order minima choose fixed branches there.
All intersections that are isolated are rational. Continuity follows
from finite maxima and minima of continuous functions. Nondecreasing
edge costs give nondecreasing cut maxima and order minima. Each active
branch has a nonnegative integral slope. A fixed tie rule chooses an
order on every open interval. Taking a minimum may destroy convexity.

Subtract the uniform fixed-order and optimal expansions. The supremum
and then the minimum preserve a uniform bounded error, proving
\eqref{eq:v26-regret-law}, even when the chosen order depends on
$\theta$. A continuous piecewise affine function on a compact interval
attains its maximum at a subdivision endpoint, proving the last claim.
These statements classify leading phases; bounded terms can still
resolve ties between orders at a transition.
\end{proof}

\subsection{The entire phase diagram of the three-edge experiment}
The star of Corollary~\ref{cor:v25-star} has leading edge weights
\[
       w_1(s)=7s,\qquad w_2(s)=\max\{6s,9s-3\},\qquad w_3(s)=s.
\]
The second edge has contact orders zero through index six, followed by
one, two and three. No new detector or adjustable formal scale is
introduced here.

\begin{corollary}[Four phases and uniform-resolution loss]
\label{cor:v26-starphase}
For this same $24$-trial experiment the leading optimal bit coefficient is
\begin{equation}\label{eq:v26-starphase}
 W(s)=\begin{cases}
 7s,&0<s\le1,\\
 10s-3,&1\le s\le3/2,\\
 8s,&3/2\le s\le3,\\
 9s-3,&s\ge3.
 \end{cases}
\end{equation}
It is not convex. An optimal leading order isolates edge $1$ before
$s=3/2$ and edge $2$ after $s=3/2$. For either bit criterion,
\begin{equation}\label{eq:v26-starinterval}
 \min_\pi\sup_{1/2\le s\le2}
 \{B_\pi^*(\theta^{2s},\theta)-B_G^*(\theta^{2s},\theta)\}
                   =\log_2(1/\theta)+O(1).
\end{equation}
In contrast, restricting the supremum to $s\in\{1/2,2\}$ gives
$\tfrac12\log_2(1/\theta)+O(1)$.
\end{corollary}
\begin{proof}
Every star order partitions its leaves before and after the centre.
Isolating edge $1$ has cost
$W_1(s)=7s$ for $s\le1$ and $W_1(s)=10s-3$ for $s\ge1$.
Isolating edge $2$ has cost $W_2(s)=8s$ for $s\le3$ and
$W_2(s)=9s-3$ for $s\ge3$. The other two partition types are
pointwise no better: the optimal partition isolates a largest weight,
as proved in Corollary~\ref{cor:v25-star}, and edge $3$ is never a
largest weight for $s>0$. Taking $\min\{W_1,W_2\}$ gives
\eqref{eq:v26-starphase}. Its slope drops from $10$ to $8$ at
$s=3/2$, proving nonconvexity.

On $[1/2,2]$, the largest value of $W_1-W$ is one, attained at $s=2$.
For $W_2-W$ it is also one, attained at $s=1$. Isolating edge $3$
or taking an empty side has excess greater than one already at
$s=1/2$: their costs there are $13/2$ and $7$, whereas $W=7/2$.
Thus $\mathfrak g_{[1/2,2]}=1$. At just the two endpoints, the two
candidate maximal excesses are respectively one and $1/2$, while the
other partitions remain worse. Apply Theorem~\ref{thm:v26-phases}
with its simultaneous bounded errors. The intermediate resolution,
not either endpoint, is essential for \eqref{eq:v26-starinterval}.
\end{proof}
